%% ma: L11-B14-0018
%% lop: 11
%% chuong: 4
%% bai: 14
%% dang: 11.14.1
%% loai: TN4
%% muc: TH
%% dap_an: "B"
%% nguon: "tai-lieu-hoc-tap-mon-toan-lop-11-theo-chuong-trinh-giao-duc-pho-thong-2018.pdf (Lớp Toán Thầy Duy, Phạm Lê Duy, THPT Chu Văn An), Chương 4 Bài 5, Phần 1, Câu 13 (tr.529)"
%% kien_thuc: "Lớp 11 Bài 14 (phép chiếu song song)"
%% trang_thai: chinh-thuc
%% ghi_chu: "Đã duyệt 03/11/2026."
\begin{ex}
Hình biểu diễn của hình hộp $ABCD.A^{\prime}B^{\prime}C^{\prime}D^{\prime}$ lên $(BC^{\prime}D)$ theo phương $AB$ là
\choice{Một tam giác.}{\True Một hình bình hành.}{Một ngũ giác.}{Một lục giác.}
\loigiai{\sloppy Phương $AB\parallel DC\parallel D^{\prime}C^{\prime}\parallel A^{\prime}B^{\prime}$. Lên $(BC^{\prime}D)$: $A,B\mapsto B$; $C\mapsto D$ (vì $CD\parallel AB$ và $D\in(BC^{\prime}D)$); $D^{\prime}\mapsto C^{\prime}$; $D\mapsto D$, $C^{\prime}\mapsto C^{\prime}$; $A^{\prime},B^{\prime}\mapsto B_1$ với $B_1$ thỏa $\overrightarrow{BB_1}=\overrightarrow{DC^{\prime}}$. Hình chiếu là hình bình hành $BDC^{\prime}B_1$.}
\end{ex}
